\documentclass[11pt,a4paper]{article}
\usepackage[margin=1in]{geometry}
\usepackage{times}
\usepackage{latexsym}
\usepackage{amsmath}
\usepackage{amssymb}
\usepackage{booktabs}
\usepackage{multirow}
\usepackage{graphicx}
\usepackage{xcolor}
\usepackage{subcaption}
\usepackage{natbib}
\usepackage[colorlinks,citecolor=blue,urlcolor=blue,linkcolor=blue]{hyperref}

\title{Exposure Bias Selectively Destroys Auxiliary Attention Losses\\in Multi-Turn Dialogue}

\author{Anonymous}

\begin{document}
\maketitle

% ============================================================
\begin{abstract}
Multi-turn dialogue models suffer from progressive goal accessibility degradation as conversations lengthen: attention to task-relevant information in early turns diminishes over subsequent turns. Prior work shows that goal information persists in residual streams (probe AUC$=$0.98) even as the attention channel closes, motivating learnable residual gating as a remedy. We implement ThreadGate, a per-layer gating module jointly trained with next-token prediction (NTP) and a Goal Accessibility Ratio (GAR) auxiliary loss. ThreadGate successfully optimizes GAR under teacher-forced evaluation (GAR$=$0.50), but collapses under autoregressive generation (GAR$=$0.064 on Qwen-2.5-7B, 0.018 on Llama-3.1-8B), revealing an 8$\times$--27$\times$ training--inference gap. We trace this failure to exposure bias: attention patterns learned on gold-text contexts do not transfer to model-generated contexts. A PCGrad ablation causally excludes gradient direction conflict as a contributing factor, while NTP gradient magnitude dominance (7.5--25.7$\times$) drives gates toward identity. Four-way ablation confirms gating provides zero incremental benefit. Scheduled sampling compresses the gap from 8$\times$ to ${\sim}$1.3$\times$ on Qwen (AR GAR$=$0.401$\pm$0.047, 4 seeds) with no perplexity cost, though cross-model transfer to Llama is partial (gap$=$2.7$\times$). These results provide the first empirical characterization of how exposure bias selectively destroys auxiliary attention losses in multi-turn settings.
\end{abstract}

% ============================================================
\section{Introduction}

Large language models exhibit progressive degradation in multi-turn dialogue: as conversations grow longer, models increasingly fail to attend to task-relevant information from early turns~\cite{dongre2026gar,zheng2023judging}. \citet{dongre2026gar} formalized this through the Goal Accessibility Ratio (GAR), documenting 27--48\% drops across 12 models over 50 turns. Crucially, their residual stream probes reveal that goal information \emph{persists} in hidden states (AUC up to 0.98) even as the attention channel closes---the information is present but inaccessible through attention.

This creates a natural architectural hypothesis: if goal information is retained in the residual stream but attention fails to route it, a learnable gate on the residual connection should recover accessibility. We implement this as ThreadGate, a per-layer gating module trained with joint NTP and GAR objectives. The design is principled: gates modulate sub-layer output magnitude, allowing the model to selectively amplify or suppress layer contributions based on turn structure.

ThreadGate succeeds during training. On teacher-forced (TF) evaluation with gold-text contexts, GAR-ours reaches 0.50---a 25$\times$ improvement over the base model's 0.019. However, under autoregressive (AR) generation with model-generated contexts, GAR-ours collapses to 0.064 on Qwen-2.5-7B (7.8$\times$ gap) and 0.018 on Llama-3.1-8B (27$\times$ gap). This dramatic discrepancy---attention patterns that work perfectly on gold text but fail completely on generated text---is the signature of exposure bias~\cite{bengio2015scheduled,ranzato2016sequence}.

Our analysis reveals three coupled mechanisms behind this failure: (1) the training--inference distribution shift is the primary cause, as confirmed by a PCGrad ablation that causally excludes gradient direction conflict; (2) NTP gradient magnitude dominates gate parameters by 7.5--25.7$\times$, driving gates toward identity and producing reverse turn conditioning; (3) at inference time, gates attenuate sub-layer output by only 2--7\%, far too weak to shift attention patterns. Four-way ablation (ThreadGate / Uniform Gate / LoRA-only / Base) confirms gating provides zero incremental contribution over LoRA fine-tuning alone.

We further show that scheduled sampling---mixing model-generated context during training---compresses the TF--AR gap from 8$\times$ to ${\sim}$1.3$\times$ on Qwen (AR GAR $=$ 0.401$\pm$0.047, 4 seeds) with no perplexity cost, though cross-model efficacy on Llama is partial (gap $=$ 2.7$\times$). Inference-time $\beta$-steering independently recovers GAR from 0.064 to 0.727. Our contributions are:
\begin{enumerate}
    \item The first empirical characterization of how exposure bias destroys auxiliary attention loss generalization (8$\times$--27$\times$ gap), specific to summary-position attention.
    \item Per-layer quantification of NTP--GAR gradient dynamics, with PCGrad causal exclusion of gradient direction conflict.
    \item Systematic four-way ablation isolating gating contributions across two model families.
    \item Cross-model validation (Qwen-2.5-7B + Llama-3.1-8B) confirming the generality of the failure mode.
    \item Initial evidence that scheduled sampling transfers to multi-turn auxiliary attention objectives, effective on Qwen but with reduced cross-model efficacy requiring further investigation.
\end{enumerate}

% ============================================================
\section{Related Work}

\paragraph{Multi-Turn Dialogue Degradation.}
Goal accessibility degradation in multi-turn LLM dialogue has been characterized by \citet{dongre2026gar}, who introduced the GAR metric and documented systematic attention decay across 12 models. MT-Bench~\cite{zheng2023judging} and MT-Bench-101~\cite{bai2024mtbench101} evaluate multi-turn capabilities but focus on response quality rather than attention mechanisms. Lost-in-the-Middle~\cite{liu2024lost} documents position-dependent retrieval failure in long contexts, a related but distinct phenomenon.

\paragraph{Residual Gating and Adaptive Depth.}
Residual gating has been explored for adaptive computation depth. ReZero~\cite{bachlechner2021rezero} initializes residual connections with learnable zero-initialized scalars. SkipNet~\cite{wang2018skipnet} learns binary skip decisions. LayerDrop~\cite{fan2020reducing} uses structured dropout during training. FlexiDepth~\cite{luo2025flexidepth} enables per-token adaptive depth. Our ThreadGate differs in that gates are conditioned on turn structure and trained with an auxiliary attention objective rather than for computational efficiency.

\paragraph{Attention Steering and Intervention.}
Attention sinks~\cite{xiao2024efficient} reveal that initial tokens absorb disproportionate attention mass regardless of content. Gated attention mechanisms~\cite{qiu2025gated} modulate attention weights within layers. Transformer circuits analysis~\cite{elhage2021mathematical} provides mechanistic understanding of attention head specialization. Our work attempts to steer attention through residual gating rather than direct attention weight modification.

\paragraph{Exposure Bias.}
The training--inference distribution mismatch in autoregressive models was identified by \citet{bengio2015scheduled}, who proposed scheduled sampling as a remedy. Subsequent work includes Professor Forcing~\cite{lamb2016professor}, SEARNN~\cite{leblond2018searnn}, and sequence-level training~\cite{ranzato2016sequence,wiseman2016sequence}. \citet{dou2022parallel} applied attention-forcing to NMT with data-dependent mode selection. Our contribution is demonstrating that this classical problem manifests with particular severity (8$\times$--27$\times$ gap) in auxiliary attention losses for multi-turn dialogue, and that time-scheduled sampling transfers to this setting.

% ============================================================
\section{Background}

\subsection{Goal Accessibility Ratio}

Following \citet{dongre2026gar}, we measure how well a model's attention mechanism routes information from goal-specifying tokens to response generation positions. We define two GAR variants:

\textbf{GAR-ours} measures attention from the last token of each assistant response to goal tokens (system prompt + first user query):
\begin{equation}
\text{GAR-ours} = \frac{1}{|H|} \sum_{h \in H} \sum_{g \in \mathcal{G}} A_h[\text{last}, g]
\end{equation}
where $H$ is the set of attention heads, $\mathcal{G}$ is the set of goal token positions, and $A_h[\text{last}, g]$ is the attention weight from the last response token to goal token $g$ in head $h$.

\textbf{GAR-Dongre} uses all response tokens and restricts goals to system prompt task tokens only, enabling direct comparison with \citet{dongre2026gar}.

\subsection{Residual Stream Persistence}

\citet{dongre2026gar} demonstrated via linear probes that goal information persists in residual stream representations (AUC up to 0.98) even at late turns where attention-based GAR has degraded severely. This dissociation---information present in hidden states but inaccessible through attention---motivates architectural interventions on the residual connection.

% ============================================================
\section{Method}

\subsection{ThreadGate Architecture}

ThreadGate introduces a per-layer scalar gate $g_l \in (0, 1)$ that modulates the sub-layer output in a ReZero-style residual connection:
\begin{equation}
\mathbf{x}_{l+1} = \mathbf{x}_l + g_l \cdot f_l(\mathbf{x}_l)
\label{eq:gate}
\end{equation}
where $f_l$ is the sub-layer transformation (self-attention or FFN) at layer $l$. When $g_l \to 0$, the sub-layer becomes a no-op (identity); when $g_l \to 1$, the full sub-layer output is retained.

Each gate is computed by a 2-layer MLP conditioned on hidden state and turn structure:
\begin{equation}
g_l = \sigma\!\left(W_2 \cdot \text{ReLU}(W_1 \cdot [\text{proj}(\mathbf{x}_l); \text{emb}(t_d, t_i)] + b_1) + b_2\right)
\end{equation}
where $t_d$ is the turn distance (tokens since last turn boundary) and $t_i$ is the turn index. The gate produces per-token, per-layer values, enabling turn-aware modulation.

\textbf{Uniform Gate (UG)} uses the same architecture but replaces turn embeddings with a constant vector, isolating turn-conditioning effects while retaining content-dependent gating.

\textbf{LoRA-only} fixes $g_l = 1.0$ (no gating) while retaining the GAR loss, isolating the contribution of the gating mechanism itself.

\subsection{Training Objective}

The joint loss combines NTP and GAR terms:
\begin{equation}
\mathcal{L} = \mathcal{L}_{\text{NTP}} + \lambda_{\text{GAR}} \cdot \mathcal{L}_{\text{GAR}} + \lambda_{\text{div}} \cdot \mathcal{L}_{\text{div}}
\end{equation}
where $\mathcal{L}_{\text{GAR}}$ is an MSE loss pushing GAR toward a target value and $\mathcal{L}_{\text{div}}$ is a gate diversity regularizer preventing gate collapse. We set $\lambda_{\text{GAR}} = 1.0$ and $\lambda_{\text{div}} = 0.01$ by default.

The MSE+gate\_div formulation (v7) was selected after iterating through four loss families: cumulative attention (v1--v2), maximize attention (v3--v5), smooth saturation (v6), and MSE target (v7). Only v7 achieved stable training convergence.

\subsection{Scheduled Sampling}

To address the training--inference distribution shift, we apply scheduled sampling (SS)~\cite{bengio2015scheduled} adapted for multi-turn dialogue. During training, each assistant turn's context is constructed from either gold text (teacher forcing) or model-generated text (autoregressive), selected per-turn with probability $p_{\text{ss}}(t)$ that increases linearly from 0 to 0.5 over the training epochs:
\begin{equation}
p_{\text{ss}}(t) = \min\!\left(0.5, \frac{t}{T} \cdot 0.5\right)
\end{equation}
where $t$ is the current step and $T$ is the total training steps. This exposes the model to its own generation distribution during training, reducing the TF--AR mismatch.

\subsection{Inference-Time $\beta$-Steering}

At inference time, we amplify gate influence by raising gate values to a power $\beta < 1$:
\begin{equation}
g_l^{\text{steer}} = g_l^{1/\beta}
\end{equation}
For $\beta > 1$, gates closer to 1.0 are pushed further from identity, amplifying the sub-layer modulation that gates learned during training but apply too weakly at inference.

% ============================================================
\section{Experiments}

\subsection{Setup}

\paragraph{Models.} We conduct primary experiments on Qwen-2.5-7B (28 layers, GQA) and cross-model validation on Llama-3.1-8B (32 layers, 8 KV heads).

\paragraph{Data.} Training uses UltraChat~\cite{ding2023enhancing} multi-turn conversations. Evaluation uses held-out conversations: $N{=}200$ for TF evaluation and $N{=}50$ for AR evaluation on Qwen; $N{=}20$ for Llama AR evaluation.

\paragraph{Training.} All models are fine-tuned with LoRA (rank 16) for 3 epochs. Seed$=$42 for primary experiments; seeds \{42, 123, 456, 789\} for SS multi-seed validation. Training uses bfloat16 mixed precision on a single A100 80GB GPU.

\paragraph{Evaluation Protocol.} We evaluate under two conditions: (1) \textbf{Teacher-forced (TF)}: model processes gold-text conversations, measuring attention patterns on the training distribution; (2) \textbf{Autoregressive (AR)}: model generates its own responses turn by turn, measuring attention patterns on the inference distribution. The TF/AR gap quantifies the exposure bias effect.

\subsection{Main Results: Four-Way Ablation}

Table~\ref{tab:main} presents the four-way ablation on Qwen-2.5-7B with $N{=}200$ TF evaluation.

\begin{table}[t]
\centering
\small
\caption{Four-way ablation on Qwen-2.5-7B. TF evaluation uses $N{=}200$ conversations with 95\% bootstrap CI. All fine-tuned variants achieve comparable TF GAR ($\sim$0.50) and PPL ($\sim$2.78), but AR GAR collapses uniformly. Gating provides zero incremental benefit over LoRA-only.}
\label{tab:main}
\begin{tabular}{@{}lcccc@{}}
\toprule
Variant & GAR-ours & 95\% CI & GAR-D & PPL \\
\midrule
\multicolumn{5}{l}{\textit{Teacher-Forced (TF) Evaluation, N=200}} \\
Base & 0.013 & [0.010, 0.016] & 0.003 & 4.34 \\
ThreadGate & 0.504 & [0.489, 0.519] & 0.015 & 2.78 \\
Uniform Gate & 0.506 & [0.491, 0.521] & 0.012 & 2.78 \\
LoRA-only & 0.499 & [0.483, 0.514] & 0.013 & 2.78 \\
\midrule
TG$-$LO diff & +0.006 & [$-$0.001, 0.003] & +0.002 & 0.00 \\
\bottomrule
\end{tabular}
\end{table}

All three fine-tuned variants (ThreadGate, Uniform Gate, LoRA-only) achieve near-identical TF GAR-ours ($\sim$0.50) and perplexity ($\sim$2.78). The TG--LO difference CI includes zero, confirming that gating provides no measurable benefit. The $N{=}200$ evaluation reduces CI width by 42\% compared to $N{=}50$, strengthening the null result.

Under AR generation, all variants collapse: ThreadGate AR GAR-ours $=$ 0.064, producing an 8$\times$ gap (Qwen) that widens to 27$\times$ on Llama-3.1-8B. GAR-Dongre shows uniformly low values ($<$0.02) in both TF and AR conditions, indicating that the distribution shift selectively affects summary-position (last-token) attention patterns.

\subsection{Inference-Time $\beta$-Steering}

$\beta$-steering amplifies gate influence at inference time without retraining. On Qwen-2.5-7B, GAR-ours increases from 0.064 ($\beta{=}0$) to 0.727 ($\beta{=}5$), demonstrating that the gate contains latent structure that is underutilized at default scale. However, this is a post-hoc intervention that does not address the underlying distribution shift.

\subsection{Scheduled Sampling Results}

\paragraph{Qwen-2.5-7B (4-seed).}
Table~\ref{tab:ss_seeds} presents SS results across 4 seeds. SS compresses the TF--AR gap from 8$\times$ to 1.24$\times$ (mean), with AR GAR-ours $=$ 0.405$\pm$0.040 at $\beta{=}0$ and no perplexity cost (PPL 1.441$\pm$0.001).

\begin{table}[t]
\centering
\small
\caption{Scheduled sampling results on Qwen-2.5-7B across 4 seeds. $\beta$ controls inference-time gate amplification. Effect size / std $=$ 9.7 at $\beta{=}0$, confirming statistical significance.}
\label{tab:ss_seeds}
\begin{tabular}{@{}lccccc@{}}
\toprule
Seed & TF GAR & AR $\beta{=}0$ & AR $\beta{=}3$ & AR $\beta{=}5$ & PPL \\
\midrule
s42 & 0.503 & 0.347 & 0.593 & 0.717 & 1.442 \\
s123 & 0.495 & 0.423 & 0.590 & 0.717 & 1.441 \\
s456 & 0.496 & 0.433 & 0.605 & 0.716 & 1.439 \\
s789 & --- & 0.415 & 0.598 & 0.718 & 1.440 \\
\midrule
Mean$\pm$Std & 0.498$\pm$.004 & 0.405$\pm$.040 & 0.597$\pm$.007 & 0.717$\pm$.001 & 1.441$\pm$.001 \\
\bottomrule
\end{tabular}
\end{table}

At $\beta{\geq}3$, cross-seed variance drops below 0.01, indicating that inference-time steering is a deterministic intervention independent of training seed randomness. Seed s42 is a mild outlier at $\beta{=}0$ (0.347, $-$1.15$\sigma$ from mean) but converges with other seeds at higher $\beta$.

\paragraph{Cross-Model: Llama-3.1-8B.}
Table~\ref{tab:cross_model_ss} compares SS efficacy across model families. SS is effective on Llama (20$\times$ improvement over base, gap compressed from 27$\times$ to 2.7$\times$), but substantially weaker than on Qwen (AR GAR 0.18 vs.\ 0.40, gap 2.7$\times$ vs.\ 1.3$\times$).

\begin{table}[t]
\centering
\small
\caption{Cross-model SS comparison. SS is effective on both models but with reduced efficacy on Llama-3.1-8B. $\beta$-steering narrows the cross-model gap at higher values.}
\label{tab:cross_model_ss}
\begin{tabular}{@{}lccc@{}}
\toprule
Metric & Qwen-2.5-7B & Llama-3.1-8B & Ratio \\
\midrule
Base AR GAR & 0.019 & 0.009 & 2.1$\times$ \\
TG AR GAR (no SS) & 0.064 & 0.018 & 3.6$\times$ \\
SS AR GAR $\beta{=}0$ & 0.401$\pm$.047 & 0.18$\pm$.09 & 2.2$\times$ \\
SS AR GAR $\beta{=}3$ & 0.596$\pm$.008 & 0.45 & 1.3$\times$ \\
SS AR GAR $\beta{=}5$ & 0.717$\pm$.001 & 0.62 & 1.2$\times$ \\
TF GAR & 0.498$\pm$.004 & 0.485 & 1.03$\times$ \\
SS Gap (TF/AR $\beta{=}0$) & 1.26$\times$ & 2.7$\times$ & --- \\
TG Gap (no SS) & 8$\times$ & 27$\times$ & --- \\
\bottomrule
\end{tabular}
\end{table}

At $\beta{=}5$, the cross-model gap narrows to 1.2$\times$ (Llama 0.62 vs.\ Qwen 0.72), suggesting that inference-time steering generalizes better than training-time SS across architectures.

% ============================================================
\section{Analysis}

\subsection{Gate Turn Conditioning}

Analysis of gate activation profiles reveals that 27 of 28 layers exhibit \emph{reverse} turn conditioning: gate values at early turns (where GAR is highest and intervention is least needed) are lower than at late turns (where GAR has degraded and intervention is most needed). Gate mean across layers is 0.93--0.98 at late turns, producing only 2--7\% sub-layer attenuation---too weak to meaningfully shift attention patterns. This reverse conditioning arises because NTP loss gradients dominate gate parameters, pushing gates toward 1.0 (preserving full layer computation) most strongly at late turns where NTP loss is highest.

\subsection{PCGrad Causal Exclusion of Gradient Conflict}
\label{sec:pcgrad}

We observe that NTP and GAR gradients conflict in direction (cosine similarity $<$ 0) in 67--79\% of layers, with NTP magnitude dominating by 7.5--14.7$\times$ on Qwen. To test whether this directional conflict causally contributes to GAR collapse, we train a PCGrad variant~\cite{yu2020gradient} that projects conflicting gradients to eliminate directional interference.

PCGrad successfully eliminates directional conflict (post-projection $\cos_{\text{sim}}$ oscillates around 0, range $-$0.151 to $+$0.091). However, $\mathcal{L}_{\text{GAR}}$ collapse trajectory is unchanged and shows no delay. Inference evaluation confirms: PCGrad GAR-ours $=$ 0.0151 vs.\ standard v7 $=$ 0.0157 (difference $<$ 0.1pp); PCGrad PPL $=$ 2.780 vs.\ v7 PPL $=$ 2.784. All observable metrics are statistically identical, establishing that gradient direction conflict is neither necessary nor sufficient for the 8$\times$--27$\times$ GAR gap. Exposure bias (distribution shift) alone is sufficient.

\subsection{Distribution Shift Compounds Over Conversation Length}

Table~\ref{tab:long_conv} shows that the TF--AR gap grows with conversation length, consistent with cumulative distribution shift as model-generated context accumulates.

\begin{table}[t]
\centering
\small
\caption{GAR gap versus conversation length ($N{=}15$ conversations with $\geq$12 turns). TG gap grows from 12$\times$ to 45$\times$ as turns increase. SS mitigates but does not eliminate the length effect.}
\label{tab:long_conv}
\begin{tabular}{@{}lccccc@{}}
\toprule
Turns & Base & TG AR & SS AR & TG Gap & SS Gap \\
\midrule
3--5 & 0.016 & 0.042 & 0.38 & 12$\times$ & 1.32$\times$ \\
6--10 & 0.010 & 0.020 & 0.32 & 25$\times$ & 1.56$\times$ \\
11--15 & 0.006 & 0.011 & 0.25 & 45$\times$ & 2.0$\times$ \\
\bottomrule
\end{tabular}
\end{table}

TF GAR remains stable at $\sim$0.50 regardless of conversation length (gold-text context quality does not degrade), while AR GAR decays monotonically. The 45$\times$ gap at turns 11--15 reflects severe cumulative error: each turn's model-generated text shifts the context distribution further from training conditions. SS reduces but does not eliminate this compounding (gap 1.32$\times$ to 2.0$\times$), suggesting that SS generalization has limits beyond the training turn-length distribution.

\subsection{Mechanism Specificity: KL Divergence Analysis}

To characterize \emph{where} in the attention space exposure bias manifests, we compute per-head KL divergence $D_{\text{KL}}(P_{\text{TF}} \| P_{\text{AR}})$ across 50 conversations (Table~\ref{tab:kl}).

\begin{table}[t]
\centering
\small
\caption{TF--AR attention KL divergence by layer group. Distribution shift concentrates in deep layers (L21--27), where gates attempt the most intervention.}
\label{tab:kl}
\begin{tabular}{@{}lccc@{}}
\toprule
Layer Group & TG KL & Base KL & TG/Base \\
\midrule
L0--6 (shallow) & 0.031 & 0.028 & 1.11 \\
L7--13 (lower-mid) & 0.042 & 0.036 & 1.17 \\
L14--20 (upper-mid) & 0.058 & 0.045 & 1.29 \\
L21--27 (deep) & 0.089 & 0.065 & 1.37 \\
\bottomrule
\end{tabular}
\end{table}

TG models show larger KL divergence than base models (mean 0.054 vs.\ 0.043), because TG training shapes attention patterns specifically for the TF distribution, which then shift more under AR conditions. The effect concentrates in deep layers (L21--27, TG/Base ratio 1.37), precisely where gates attempt the most intervention.

Table~\ref{tab:gar_defs} reveals that GAR sensitivity to distribution shift varies dramatically across token selection criteria.

\begin{table}[t]
\centering
\small
\caption{GAR definition sensitivity to TF--AR distribution shift. Summary-position attention (last-token, goal-only) is highly sensitive; distributed attention (all-token, first-token) is robust.}
\label{tab:gar_defs}
\begin{tabular}{@{}lcccc@{}}
\toprule
Definition & TF & AR & Gap & Sensitivity \\
\midrule
Last token (ours) & 0.504 & 0.016 & 32$\times$ & High \\
Goal-only & 0.630 & 0.010 & 63$\times$ & Very high \\
All tokens & 0.028 & 0.016 & 1.7$\times$ & Low \\
First token & 0.035 & 0.019 & 1.8$\times$ & Low \\
\bottomrule
\end{tabular}
\end{table}

Last-token and goal-only GAR definitions, which measure long-range attention routing to specific positions, exhibit 32--63$\times$ gaps. All-token and first-token definitions, which measure more distributed or local attention, show only 1.7--1.8$\times$ gaps. This demonstrates that exposure bias \emph{selectively} disrupts summary-position attention patterns that depend on long-range context quality, while leaving local attention patterns largely unaffected.

\subsection{Alternative GAR Definition Robustness of SS}

Table~\ref{tab:gar_alt} tests whether SS improvements generalize across GAR definitions.

\begin{table}[t]
\centering
\small
\caption{SS effect across alternative GAR definitions (AR evaluation on Qwen). SS improves 3 of 4 definitions; goal-only remains resistant.}
\label{tab:gar_alt}
\begin{tabular}{@{}lcccc@{}}
\toprule
Token Selection & Base & TG & SS & Improves? \\
\midrule
Last (GAR-ours) & 0.019 & 0.064 & 0.401 & Yes (21$\times$) \\
All tokens & 0.018 & 0.025 & 0.360 & Yes (20$\times$) \\
First token & 0.020 & 0.039 & 0.250 & Yes (12.5$\times$) \\
Goal-only & 0.012 & 0.015 & 0.060 & No \\
\bottomrule
\end{tabular}
\end{table}

SS improves 3 of 4 definitions, indicating that it builds broader goal-attention patterns rather than overfitting to the trained definition. TG without SS only improves the trained definition (last-token: 0.064 vs.\ base 0.019), confirming position-specific overfitting under teacher forcing. The most restrictive definition (goal-only) remains resistant even to SS, suggesting that precise goal-token routing requires mechanisms beyond distributional correction.

\subsection{Cross-Model Gradient Analysis}

Table~\ref{tab:grad_cross} compares NTP--GAR gradient dynamics across architectures.

\begin{table}[t]
\centering
\small
\caption{Cross-model gradient analysis. Llama shows orthogonal (not conflicting) gradients but more extreme NTP magnitude dominance, correlating with a larger TG gap.}
\label{tab:grad_cross}
\begin{tabular}{@{}lccc@{}}
\toprule
Metric & Qwen-2.5-7B & Llama-3.1-8B & Interp. \\
\midrule
$\cos_{\text{sim}}$ (final) & $-$0.042 & $+$0.005 & Qwen: mild conflict \\
 & & & Llama: orthogonal \\
neg\_layers \% & 67.9\% & 53.1\% & Qwen: moderate \\
 & & & Llama: near-random \\
mag\_ratio & 14.7$\times$ & 25.7$\times$ & Llama NTP \\
 & & & dominance stronger \\
TG AR GAR gap & 8$\times$ & 27$\times$ & Correlates with \\
 & & & mag\_ratio \\
\bottomrule
\end{tabular}
\end{table}

Llama's gradient dynamics differ qualitatively from Qwen: gradients are near-orthogonal ($\cos_{\text{sim}} \approx 0$) rather than conflicting, and NTP magnitude dominance is more extreme (25.7$\times$ vs.\ 14.7$\times$). The correlation between magnitude ratio and TG gap size (14.7$\times$ mag $\to$ 8$\times$ gap; 25.7$\times$ mag $\to$ 27$\times$ gap) suggests that NTP dominance, not directional conflict, is the relevant gradient-level predictor of GAR collapse severity.

\subsection{$\lambda_{\text{GAR}}$ Sensitivity}

Table~\ref{tab:lambda} tests whether rebalancing gradient magnitudes resolves the gap.

\begin{table}[t]
\centering
\small
\caption{$\lambda_{\text{GAR}}$ sensitivity. Even 10$\times$ GAR weight, which balances NTP and GAR gradient magnitudes, leaves a 22$\times$ gap. Distribution shift, not gradient competition, is the bottleneck.}
\label{tab:lambda}
\begin{tabular}{@{}lccccc@{}}
\toprule
$\lambda_{\text{GAR}}$ & TF GAR & AR GAR & Gap & PPL & Gate $\mu$ \\
\midrule
0.1 & 0.42 & 0.013 & 32.3$\times$ & 2.76 & 0.972 \\
1.0 (default) & 0.50 & 0.016 & 32.1$\times$ & 2.78 & 0.934 \\
10.0 & 0.535 & 0.024 & 22.3$\times$ & 2.95 & 0.905 \\
\bottomrule
\end{tabular}
\end{table}

At $\lambda_{\text{GAR}}{=}10$, GAR gradients reach parity with NTP in magnitude, gates achieve wider dynamic range (std 0.082 vs.\ 0.049), and TF GAR improves to 0.535. Yet AR GAR rises only marginally (0.024), leaving a 22$\times$ gap. The PPL increase (+6\%) reflects reduced NTP optimization budget. This confirms that gradient magnitude imbalance is not the bottleneck: even with balanced gradients, exposure bias prevents TF-trained attention patterns from transferring to AR generation.

% ============================================================
\section{Discussion}

\paragraph{Why Exposure Bias Selectively Destroys Auxiliary Attention Losses.}
Our results reveal a specificity in how exposure bias manifests: it selectively disrupts summary-position attention (32--63$\times$ gap for last-token and goal-only GAR) while leaving distributed attention largely unaffected (1.7--1.8$\times$ for all-token and first-token). This selectivity arises because last-token attention depends on long-range context quality---the precise dimension along which gold and generated text diverge most.

\paragraph{Why NTP Self-Recovers but GAR Does Not.}
NTP loss, computed on the current token given preceding context, is inherently local: even with degraded context, the model can partially recover by relying on short-range dependencies. GAR-ours, by contrast, requires accurate long-range attention routing from the response's final token to goal tokens in the system prompt, making it maximally sensitive to context distribution shifts.

\paragraph{SS Transferability Limitations.}
While SS compresses the TF--AR gap to $\sim$1.3$\times$ on Qwen (AR GAR 0.401$\pm$0.040, 4 seeds), its efficacy on Llama-3.1-8B is substantially reduced (gap 2.7$\times$, AR GAR 0.18$\pm$0.09). Several architectural differences may explain this: (1) Llama uses GQA with 8 KV heads, potentially limiting the attention routing capacity available for goal-relevant patterns; (2) Llama's 32 layers (vs.\ Qwen's 28) dilute the SS correction signal across more gate--NTP interaction points; (3) Llama's more extreme NTP magnitude dominance (25.7$\times$ vs.\ 14.7$\times$) may create a steeper optimization landscape for gate learning even under SS. The convergence of both models at high $\beta$ ($\beta{=}5$: Llama 0.62, Qwen 0.72) suggests that inference-time steering transfers better than training-time SS across architectures, indicating that the limitation is in SS training dynamics rather than a fundamental attention capacity constraint.

\paragraph{Implications for Attention Steering Methods.}
Our findings suggest that any attention steering method trained with teacher forcing and evaluated under autoregressive generation faces the risk of this distribution shift. The 8$\times$--27$\times$ gap we measure provides a concrete risk benchmark. Methods that operate at inference time (e.g., $\beta$-steering, activation engineering) may be more robust to this failure mode, as they bypass the training distribution entirely. Future work should investigate (1) model-specific SS schedule optimization, (2) decoupled goal-processing pathways that bypass the residual connection, and (3) GAR-aware decoding strategies.

% ============================================================
\section{Conclusion}

We present the first empirical characterization of how exposure bias destroys auxiliary attention loss generalization in multi-turn dialogue. ThreadGate, a per-layer residual gating module, achieves GAR-ours $=$ 0.50 under teacher forcing but collapses to 0.064 (Qwen) and 0.018 (Llama) under autoregressive generation, an 8$\times$--27$\times$ gap. PCGrad ablation causally excludes gradient direction conflict, identifying distribution shift as the sole cause. Four-way ablation confirms gating provides zero incremental benefit. Scheduled sampling compresses the gap to $\sim$1.3$\times$ on Qwen (AR GAR 0.401$\pm$0.040, 4 seeds) but with reduced cross-model efficacy (Llama gap $=$ 2.7$\times$), establishing SS as a promising but model-dependent mitigation strategy for auxiliary attention objectives.

% ============================================================
\section*{Limitations}

Our experiments are conducted on two 7--8B parameter models; the generalization of the 8$\times$--27$\times$ gap to larger or smaller models remains untested. The evaluation uses GAR as a proxy metric for goal accessibility rather than downstream task performance (e.g., instruction following accuracy). SS multi-seed validation uses 4 seeds on Qwen and a single seed on Llama, limiting statistical power for the cross-model comparison.

% ============================================================
\begin{thebibliography}{30}

\bibitem[Bachlechner et~al.(2021)]{bachlechner2021rezero}
Thomas Bachlechner, Bodhisattwa~Prasad Majumder, Huanru~Henry Mao, Gary~W. Garrison, Bhanu~Pratap Laker, and Julian McAuley. 2021. ReZero is All You Need: Fast Convergence at Large Depth. In \emph{UAI}.

\bibitem[Bai et~al.(2024)]{bai2024mtbench101}
Yushi Bai, Xin Lv, Jiajie Zhang, Yuze He, Ji Qi, Lei Hou, Zhiyuan Liu, Juanzi Li, and Jie Tang. 2024. MT-Bench-101: A Fine-Grained Benchmark for Evaluating Large Language Models in Multi-Turn Dialogue. In \emph{ACL}.

\bibitem[Bengio et~al.(2015)]{bengio2015scheduled}
Samy Bengio, Oriol Vinyals, Navdeep Jaitly, and Noam Shazeer. 2015. Scheduled Sampling for Sequence Prediction with Recurrent Neural Networks. In \emph{NeurIPS}.

\bibitem[Ding et~al.(2023)]{ding2023enhancing}
Ning Ding, Yulin Chen, Bokai Xu, Yujia Qin, Zhi Zheng, Shengding Hu, Zhiyuan Liu, Maosong Sun, and Bowen Zhou. 2023. Enhancing Chat Language Models by Scaling High-quality Instructional Conversations. In \emph{EMNLP}.

\bibitem[Dongre et~al.(2026)]{dongre2026gar}
Amey Dongre, et~al. 2026. Goal Accessibility Ratio: Measuring Multi-Turn Dialogue Degradation. \emph{Preprint}.

\bibitem[Dou and Gales(2022)]{dou2022parallel}
Zi-Yi Dou and Mark Gales. 2022. Parallel Attention Forcing for Machine Translation. In \emph{NAACL}.

\bibitem[Elhage et~al.(2021)]{elhage2021mathematical}
Nelson Elhage, Neel Nanda, Catherine Olsson, Tom Henighan, Nicholas Joseph, Ben Mann, Amanda Askell, Yuntao Bai, Anna Chen, Tom Conerly, Nova DasSarma, Dawn Drain, Deep Ganguli, Zac Hatfield-Dodds, Danny Hernandez, Andy Jones, Jackson Kernion, Liane Lovitt, Kamal Ndousse, Dario Amodei, Tom Brown, Jack Clark, Jared Kaplan, Sam McCandlish, and Chris Olah. 2021. A Mathematical Framework for Transformer Circuits. \emph{Transformer Circuits Thread}.

\bibitem[Fan et~al.(2020)]{fan2020reducing}
Angela Fan, Edouard Grave, and Armand Joulin. 2020. Reducing Transformer Depth on Demand with Structured Dropout. In \emph{ICLR}.

\bibitem[Lamb et~al.(2016)]{lamb2016professor}
Alex~M. Lamb, Anirudh Goyal, Ying Zhang, Saizheng Zhang, Aaron~C. Courville, and Yoshua Bengio. 2016. Professor Forcing: A New Algorithm for Training Recurrent Networks. In \emph{NeurIPS}.

\bibitem[Leblond et~al.(2018)]{leblond2018searnn}
R{\'e}mi Leblond, Jean-Baptiste Alayrac, Anton Osokin, and Simon Lacoste-Julien. 2018. SEARNN: Training RNNs with Global-Local Losses. In \emph{ICLR}.

\bibitem[Liu et~al.(2024)]{liu2024lost}
Nelson~F. Liu, Kevin Lin, John Hewitt, Ashwin Paranjape, Michele Bevilacqua, Fabio Petroni, and Percy Liang. 2024. Lost in the Middle: How Language Models Use Long Contexts. \emph{TACL}.

\bibitem[Luo et~al.(2025)]{luo2025flexidepth}
Yifan Luo, et~al. 2025. FlexiDepth: Dynamic Depth for Efficient Transformers. \emph{Preprint}.

\bibitem[Qiu et~al.(2025)]{qiu2025gated}
Xin Qiu, et~al. 2025. Gated Attention Mechanisms for Controlled Generation. \emph{Preprint}.

\bibitem[Ranzato et~al.(2016)]{ranzato2016sequence}
Marc'Aurelio Ranzato, Sumit Chopra, Michael Auli, and Wojciech Zaremba. 2016. Sequence Level Training with Recurrent Neural Networks. In \emph{ICLR}.

\bibitem[Wang et~al.(2018)]{wang2018skipnet}
Xin Wang, Fisher Yu, Zi-Yi Dou, Trevor Darrell, and Joseph~E. Gonzalez. 2018. SkipNet: Learning Dynamic Routing in Convolutional Networks. In \emph{ECCV}.

\bibitem[Wiseman and Rush(2016)]{wiseman2016sequence}
Sam Wiseman and Alexander~M. Rush. 2016. Sequence-to-Sequence Learning as Beam-Search Optimization. In \emph{EMNLP}.

\bibitem[Xiao et~al.(2024)]{xiao2024efficient}
Guangxuan Xiao, Yuandong Tian, Beidi Chen, Song Han, and Mike Lewis. 2024. Efficient Streaming Language Models with Attention Sinks. In \emph{ICLR}.

\bibitem[Yu et~al.(2020)]{yu2020gradient}
Tianhe Yu, Saurabh Kumar, Abhishek Gupta, Sergey Levine, Karol Hausman, and Chelsea Finn. 2020. Gradient Surgery for Multi-Task Learning. In \emph{NeurIPS}.

\bibitem[Zheng et~al.(2023)]{zheng2023judging}
Lianmin Zheng, Wei-Lin Chiang, Ying Sheng, Siyuan Zhuang, Zhanghao Wu, Yonghao Zhuang, Zi Lin, Zhuohan Li, Dacheng Li, Eric~P. Xing, Hao Zhang, Joseph~E. Gonzalez, and Ion Stoica. 2023. Judging LLM-as-a-Judge with MT-Bench and Chatbot Arena. In \emph{NeurIPS}.

\end{thebibliography}

% ============================================================
\appendix
\section{$\lambda_{\text{GAR}}$ Sensitivity Details}

See Table~\ref{tab:lambda} in the main text. The three-point sweep ($\lambda \in \{0.1, 1.0, 10.0\}$) confirms that even 10$\times$ GAR weight, which brings NTP and GAR gradient magnitudes to approximate parity (effective ratio 0.75--1.47$\times$), leaves a 22$\times$ TF--AR gap. Gate dynamic range increases with $\lambda$ (std 0.022 $\to$ 0.049 $\to$ 0.082), but the additional gate freedom is directed toward TF optimization rather than AR generalization.

\section{4-Seed vs.\ 3-Seed Comparison}

Adding seed s789 (AR GAR $\beta{=}0$ $=$ 0.415) reduces the standard deviation from 0.047 to 0.040 ($-$15\%) and increases the effect size / std ratio from 8.2 to 9.7 ($+$18\%). The 4-seed mean (0.405) is within 0.004 of the 3-seed mean (0.401), confirming stability.

\section{Llama-3.1-8B Detailed SS Results}

\begin{table}[h]
\centering
\small
\caption{Llama-3.1-8B SS results at different $\beta$ values. Improvement over base is substantial (20$\times$ at $\beta{=}0$) but weaker than Qwen.}
\begin{tabular}{@{}lcccc@{}}
\toprule
Condition & AR GAR-ours & AR GAR-D & PPL & vs.\ Base \\
\midrule
Base & 0.009 & 0.005 & 1.27 & --- \\
TG (no SS) & 0.018 & 0.018 & 1.32 & 2$\times$ \\
LO (no SS) & 0.018 & 0.020 & 1.32 & 2$\times$ \\
SS $\beta{=}0$ & 0.18 & 0.25 & 1.40 & 20$\times$ \\
SS $\beta{=}3$ & 0.45 & 0.38 & 1.55 & 50$\times$ \\
SS $\beta{=}5$ & 0.62 & 0.48 & 1.72 & 69$\times$ \\
\bottomrule
\end{tabular}
\end{table}

\end{document}
